\documentclass[12pt,oneside]{article}

\usepackage[letterpaper,margin=1in,headheight=15pt,footskip=30pt]{geometry}
\usepackage[T1]{fontenc}
\usepackage[utf8]{inputenc}
\usepackage{mathpazo}
\usepackage[scaled=0.92]{helvet}
\usepackage{courier}
\usepackage{microtype}
\usepackage{amsmath,amssymb}
\usepackage{xcolor}
\usepackage{parskip}
\usepackage{setspace}
\usepackage{enumitem}
\usepackage{titlesec}
\usepackage{fancyhdr}
\usepackage{lastpage}
\usepackage{hyperref}
\usepackage{xurl}
\usepackage{bookmark}
\usepackage[most]{tcolorbox}
\usepackage{ragged2e}
\usepackage{etoolbox}
\usepackage{needspace}
\usepackage{endnotes}

\definecolor{ManifestInk}{HTML}{202124}
\definecolor{ManifestGray}{HTML}{5F6368}
\definecolor{ManifestLight}{HTML}{F3F4F5}
\definecolor{ManifestRule}{HTML}{8A8D91}

\hypersetup{
  colorlinks=true,
  linkcolor=ManifestInk,
  urlcolor=ManifestGray,
  pdftitle={When Caring Becomes Camouflage: An Ethics of Quiet
Quitting},
  pdfauthor={De-identified working draft},
  pdfsubject={A case-anchored argument for repair, legibility, and the
right to withdraw hidden heroic labor},
  pdfkeywords={quiet quitting, ethics, work, labor, management, safety, employee voice}
}

\doublespacing
\setlength{\parindent}{0pt}
\setlength{\parskip}{0.45em}
\setcounter{secnumdepth}{-1}
\setcounter{tocdepth}{2}
\emergencystretch=3em
\clubpenalty=10000
\widowpenalty=10000
\displaywidowpenalty=10000

\titleformat{\section}
  {\Large\sffamily\bfseries\color{ManifestInk}}
  {}{0pt}{}
  [\vspace{0.2em}\color{ManifestRule}\titlerule]
\titlespacing*{\section}{0pt}{2.2em}{0.8em}

\titleformat{\subsection}
  {\large\sffamily\bfseries\color{ManifestInk}}
  {}{0pt}{}
\titlespacing*{\subsection}{0pt}{1.5em}{0.45em}

\pretocmd{\section}{\Needspace{6\baselineskip}}{}{}
\pretocmd{\subsection}{\Needspace{4\baselineskip}}{}{}

\setlist[itemize]{leftmargin=1.45em,itemsep=0.22em,topsep=0.35em}
\setlist[enumerate]{leftmargin=1.65em,itemsep=0.32em,topsep=0.4em}
\providecommand{\tightlist}{%
  \setlength{\itemsep}{0.22em}\setlength{\parskip}{0pt}}

\renewenvironment{quote}
  {\begin{tcolorbox}[
      enhanced,
      breakable,
      colback=ManifestLight,
      colframe=ManifestRule,
      boxrule=0pt,
      leftrule=3pt,
      arc=0pt,
      outer arc=0pt,
      left=10pt,right=10pt,top=8pt,bottom=8pt,
      before skip=0.8em,after skip=0.8em
    ]\itshape}
  {\end{tcolorbox}}

\pagestyle{fancy}
\fancyhf{}
\fancyhead[L]{\sffamily\footnotesize\color{ManifestGray}When Caring
Becomes Camouflage}
\fancyhead[R]{\sffamily\footnotesize\color{ManifestGray}Working draft}
\fancyfoot[C]{\sffamily\footnotesize\color{ManifestGray}\thepage\ / \pageref*{LastPage}}
\renewcommand{\headrulewidth}{0.3pt}
\renewcommand{\footrulewidth}{0pt}

\let\footnote=\endnote
\renewcommand{\notesname}{Notes}

\begin{document}
\hypersetup{pageanchor=false}

\begin{titlepage}
\thispagestyle{empty}
\vspace*{0.65in}

{\sffamily\bfseries\fontsize{26}{31}\selectfont\color{ManifestInk}
When Caring Becomes Camouflage: An Ethics of Quiet Quitting\par}

\vspace{0.5em}
{\color{ManifestRule}\rule{\textwidth}{1.2pt}}
\vspace{1.0em}

{\sffamily\fontsize{15}{20}\selectfont\color{ManifestGray}
A case-anchored argument for repair, legibility, and the right to
withdraw hidden heroic labor\par}

\vfill

\begin{tcolorbox}[
  enhanced,
  colback=ManifestLight,
  colframe=ManifestRule,
  boxrule=0pt,
  leftrule=4pt,
  arc=0pt,
  left=14pt,right=14pt,top=13pt,bottom=13pt
]
{\Large\itshape A last boundary, not a first weapon.}
\end{tcolorbox}
\vfill

{\sffamily\color{ManifestGray}August 16, 2026\par}
\vspace{0.4em}
{\sffamily\bfseries\color{ManifestGray}Working draft - not cleared for
submission\par}
\vspace{0.55in}
\end{titlepage}

\pagenumbering{arabic}
\hypersetup{pageanchor=true}

\hypertarget{abstract}{%
\subsection{Abstract}\label{abstract}}

Quiet quitting is usually framed as a contest between worker entitlement
and managerial extraction. This essay offers a narrower ethical account
grounded in lived workplace experience. A worker may stop supplying
unsustainable discretionary labor when an organization has come to
depend on that hidden subsidy, reasonable and safe repair pathways have
failed, and competent work, honesty, safety, proportionality, and
concern for other people remain fixed. I call the withdrawal
\emph{ethical de-buffering}: it makes actual capacity legible without
manufacturing failure. The essay also names the \emph{fuck-it factor}, a
phenomenological threshold at which acute constraint conflict becomes
conscious triage. That phrase is not a validated scientific construct or
an excuse. Together, the concepts distinguish an accountable last
boundary from apathy, sabotage, or negligent work, while locating the
larger repair duty with those who control shared resources and
priorities.

\textbf{Keywords:} quiet quitting; employee voice; job quality; work
organization; occupational safety

\hypertarget{when-caring-becomes-camouflage}{%
\section{When Caring Becomes
Camouflage}\label{when-caring-becomes-camouflage}}

I did not arrive at quiet quitting because I stopped caring. I arrived
there by caring in a way that made a broken arrangement look healthier
than it was.

The pattern will be familiar to many workers, although I do not claim it
is universal. A workplace has more demand than its official staffing,
equipment, or process can reliably meet. Someone notices the gap and
closes it with private effort: moving faster, skipping recovery,
absorbing a vacancy, solving problems off the clock, carrying the
emotional load, or treating every predictable rush as a personal
emergency. The work gets done. The rescue disappears into the number.
Tomorrow, the number becomes proof that the official resources were
sufficient.

Competence becomes evidence against the competent worker's request for
help.

This is the case from which my argument begins. It is not a survey of
everyone called a quiet quitter. It is a conditional ethical claim: when
an institution has come to depend on unsustainable discretionary excess,
when reasonable and safe repair pathways have failed, and when the
worker preserves competent core performance, safety, honesty,
proportionality, and regard for other people, withdrawing that excess
can be an ethical last boundary.

That practice deserves a more exact name than either ``laziness'' or
``disengagement.'' I call it \textbf{quiet quitting proper}. Quiet does
not mean secret. It means that the unowned donation stops. The worker
continues the real job, but the undefined second job - the chronic
rescue operation hidden inside the first - no longer belongs to the
employer by default.

The distinction is simple: ceasing to subsidize a system is not the same
act as manufacturing damage to it.

\hypertarget{pathways-can-close-before-a-worker-leaves}{%
\subsection{Pathways Can Close Before a Worker
Leaves}\label{pathways-can-close-before-a-worker-leaves}}

The usual response is: if the job is so bad, why not leave?

Sometimes leaving is exactly right. But the question treats a formal
option as though it were automatically a live one. A vacancy can exist
without being a viable transition for this worker at this time. Search
consumes attention, schedule control, transportation, references,
health, and tolerance for uncertainty. A new job may mean a gap in
income, lower pay, lost insurance, an incompatible schedule, or the risk
that the current employer will end the wage before the next one begins.

The material constraints are not hypothetical. In the Federal Reserve's
2025 household survey, 16 percent of adults said they had not paid all
their bills in full in the prior month, and only 63 percent said they
could cover a \$400 emergency with cash or its equivalent.\footnote{Board
  of Governors of the Federal Reserve System, \emph{Economic Well-Being
  of U.S. Households in 2025} (Washington, DC, May 2026),
  https://www.federalreserve.gov/publications/files/2025-report-economic-well-being-us-households-202605.pdf.}
A 2026 Gallup-West Health survey found that 24 percent of employed
adults who relied on an employer for their primary health coverage said
they stayed in a job they wanted to leave because they needed the
insurance.\footnote{West Health and Gallup, ``One in Four U.S. Employees
  Locked in Jobs for Health Insurance,'' July 21, 2026,
  https://news.gallup.com/poll/712166/one-four-employees-locked-jobs-health-insurance.aspx.}
Those numbers do not establish that quiet quitting is justified for any
individual. They establish that ``just leave'' often deletes the
transition problem.

Employment conditions can also consume the resources needed to exit. An
observational study of hourly service workers associated unstable
schedules with distress, poor sleep, and unhappiness.\footnote{Daniel
  Schneider and Kristen Harknett, ``Consequences of Routine
  Work-Schedule Instability for Worker Health and Well-Being,''
  \emph{American Sociological Review} 84, no. 1 (2019): 82-114,
  https://doi.org/10.1177/0003122418823184.} Quiet quitting proper can
sometimes preserve income while recovering enough time or health to
search, train, organize, or leave. It can also lead to discipline, lost
advancement, or dismissal. It is not a guaranteed bridge. It is one
possible reconfiguration of constrained agency.

Exit is not the only pathway that can close. Voice can become formally
available but causally inert. The worker reports a conflict, yet nobody
with authority chooses among the demands. A suggestion is praised and
forgotten. Excellent work is rewarded with a larger unresourced role. A
chain of command transmits commands downward while editing bad news on
the way up.

Research on employee voice and silence supports the limited point that
speaking depends partly on whether workers expect it to be safe and
effective; organizations can lose useful information when silence
appears more rational than voice.\footnote{Elizabeth W. Morrison,
  ``Employee Voice and Silence,'' \emph{Annual Review of Organizational
  Psychology and Organizational Behavior} 1 (2014): 173-197,
  https://doi.org/10.1146/annurev-orgpsych-031413-091328.} It does not
prove that every rejected proposal closes a pathway. A good-faith answer
the worker dislikes is not closure. Closure is better shown by a
repeated material pattern: nonresponse, retaliation, refusal to rank
demonstrated incompatibilities, or a nominal remedy with no operative
route.

\hypertarget{repair-first-but-not-forever}{%
\subsection{Repair First, But Not
Forever}\label{repair-first-but-not-forever}}

Quiet quitting proper should be a last boundary, not a first weapon.
That requires repair first - but repair first must not become another
demand for martyrdom.

A worker does not need institutional permission to stop making a
genuinely voluntary gift. The stronger repair obligation applies when an
employer has come to rely on an established pattern of excess, when
other people could be harmed by abrupt withdrawal, or when the worker
intends the boundary to make a structural problem visible. Even then,
duties scale with information, authority, safety, and remaining
capacity. A line worker is not required to perform an unpaid operations
study before saying, ``At this pace and staffing, I cannot complete A,
B, and C safely.''

The nearest functioning channel may be a supervisor, coworker, union
steward, safety officer, formal complaint process, or protected external
route. It need not be the nearest person on the organization chart. A
worker should not have to repeat a warning through a channel reasonably
expected to retaliate, conceal danger, or destroy evidence. Urgent
safety and legal issues can justify bypassing ordinary sequence.

Where possible, the cleanest instrument is an explicit priority request:

\begin{quote}
With the current staffing, demand, and equipment, we can complete A and
B safely, or A and C safely, but not A, B, and C at the demanded
standard. Which has priority?
\end{quote}

That question transfers the tradeoff to the level that owns the
decision. A front-line supervisor may not control labor or budgets;
their repair duty may be to validate the conflict and carry it upward
with the worker. Repair duties follow decision rights. Humanness is
symmetric. Power is not.

The worker should preserve a lawful and appropriate contemporaneous
account of conditions: staffing, demand, shortages, equipment failures,
warnings, priority decisions, and resulting noncritical backlog or
delay. This record is not proof of one unique cause. It keeps the
resource context from vanishing before rival explanations are tested.

Do not create a failure to send a message. Stop erasing evidence of a
failure that already exists.

Repair can end. Repeated nonresponse, retaliation, a refusal to own
demonstrated tradeoffs, or a plainly futile loop can close the
reasonable path. Workers do not owe infinite appeals. Closure is also
revisable. A material repair deserves reassessment, but it does not
create an automatic claim to the old level of heroics.

\hypertarget{ethical-de-buffering}{%
\subsection{Ethical De-Buffering}\label{ethical-de-buffering}}

An organization may possess more observed capacity than sustainable
capacity because a worker supplies a hidden buffer. The hidden input
includes skipped recovery, chronic vacancy coverage, off-clock
troubleshooting, exceptional flexibility, or a pace that can be survived
briefly but not repeated.

I call the proportionate withdrawal of that input \textbf{ethical
de-buffering}. The worker returns hidden heroic labor toward a
sustainably repeatable, reasonably agreed or compensated baseline while
holding competent core work, safety, and honesty fixed.

The adjective \emph{ethical} is not self-certifying. Some jobs
legitimately include initiative, seasonal surges, variable schedules, or
broad responsibility. A skilled worker's high output is not
automatically exploitation. The core role must be inferred from the
lawful agreement, written duties, professional standards, compensation,
specific reasonable instructions, and settled practice, all bounded by
safety and feasibility. Repetition alone does not convert an
unsustainable subsidy into a duty.

Ethical de-buffering can mean taking agreed breaks, ending habitual
off-clock rescue, declining indefinite vacancy coverage, requesting
authorization before accepting a materially expanded role, and allowing
noncritical backlog to remain visible. Established reliance can create a
proportionate notice or handoff duty where feasible. It cannot create
permanent ownership of the worker's health.

A drop in output after de-buffering proves only that the operation
depended on the withdrawn effort. It does not prove understaffing or
identify a single cause. Demand, process design, equipment, skill mix,
scheduling, absenteeism, and individual performance remain live
explanations. The change is a diagnostic signal, not a verdict.

The serious boundary is human harm. Noncritical delay can reveal
sustainable capacity. Foreseeable injury, contamination, medical error,
or public danger is not an acceptable signal. If no productive action
preserves the safety floor, the correct response may be to slow, stop,
refuse, transfer, shut down, or escalate. Safety is not one weight among
many when the competing objective is throughput.

This matters because hidden heroics may be the last margin in a brittle
system. Near a bottleneck, a small loss of effective capacity can
sometimes create a much larger increase in queues, coordination failure,
or delay. Queueing and resilience models make that possibility
intelligible, but they do not predict every workplace.\footnote{J. F. C.
  Kingman, ``The Single Server Queue in Heavy Traffic,''
  \emph{Proceedings of the Cambridge Philosophical Society} 57, no. 4
  (1961): 902-904, https://doi.org/10.1017/S0305004100036094; David D.
  Woods, ``The Theory of Graceful Extensibility,'' \emph{Environment
  Systems and Decisions} 38 (2018): 433-457,
  https://doi.org/10.1007/s10669-018-9708-3.} A nonlinear disruption may
reveal dependence on hidden adaptive capacity without proving what
caused the bottleneck.

The worker's body is not the slack variable in a staffing model.

\hypertarget{the-fuck-it-factor}{%
\subsection{The Fuck-It Factor}\label{the-fuck-it-factor}}

The long-term story has a short-term counterpart. A workplace can close
pathways over months and still fail inside a single rush.

\textbf{The fuck-it factor} is my term for an acute felt transition: a
person moves from trying to preserve every live requirement to judging,
under time pressure and reduced clarity, that no fully compliant action
appears reachable before a response is required.

The phrase carries the phenomenological truth of the moment: \emph{Fuck
it. I cannot preserve everything. I will try to sacrifice the least
harmful thing I am ethically permitted to sacrifice.}

This is my coinage, not an established human-factors construct,
validated scale, clinical diagnosis, legal defense, or excuse. The
action set may truly be empty, or it may merely appear empty within the
time, information, equipment, and authority available to the person. A
later investigation may find a route the worker could not see. That
difference matters for prevention and accountability.

The surrounding risk conditions are better established than the coined
threshold. Human-factors research shows that time pressure can produce
speed-accuracy tradeoffs, acute stress can impair some executive
functions, interruptions can disrupt sequences, and production pressure
or competing goals can contribute to procedural departures. The effects
vary by task, expertise, and setting.\footnote{James L. Szalma, Peter A.
  Hancock, and Sean Quinn, ``A Meta-Analysis of the Effect of Time
  Pressure on Human Performance,'' \emph{Proceedings of the Human
  Factors and Ergonomics Society Annual Meeting} 52, no. 19 (2008):
  1513-1516, https://doi.org/10.1177/154193120805201944; Grant S.
  Shields, Matthew A. Sazma, and Andrew P. Yonelinas, ``The Effects of
  Acute Stress on Core Executive Functions,'' \emph{Neuroscience \&
  Biobehavioral Reviews} 68 (2016): 651-668,
  https://doi.org/10.1016/j.neubiorev.2016.06.038; Johanna I. Westbrook
  et al., ``Association of Interruptions with an Increased Risk and
  Severity of Medication Administration Errors,'' \emph{Archives of
  Internal Medicine} 170, no. 8 (2010),
  https://doi.org/10.1001/archinternmed.2010.65; Seyed M. Hashemian and
  Konstantinos Triantis, ``Production Pressure and Its Relationship to
  Safety: A Systematic Review,'' \emph{Safety Science} 159 (2023):
  106045, https://doi.org/10.1016/j.ssci.2022.106045.}

Three events must remain distinct. An error is unintentional. A
procedural departure is intentional only in the sense that the person
knows a rule or requirement is being relaxed; it does not imply a desire
for harm. A mixed cascade begins with a deliberate shortcut and produces
an unintended result. Human-factors researchers distinguish errors from
violations, while research on workarounds shows that local adaptation
can both preserve necessary work and create hidden risk.\footnote{James
  Reason et al., ``Errors and Violations on the Roads: A Real
  Distinction?'' \emph{Ergonomics} 33, nos. 10-11 (1990): 1315-1332,
  https://doi.org/10.1080/00140139008925335; Debra S. Debono et al.,
  ``Nurses' Workarounds in Acute Healthcare Settings: A Scoping
  Review,'' \emph{BMC Health Services Research} 13 (2013): 175,
  https://doi.org/10.1186/1472-6963-13-175.}

Deliberate triage may reduce total harm when full compliance is
impossible. It may also open a failure path the person cannot see. We do
not yet know whether the felt transition I name adds risk beyond the
severe conditions that trigger it. That question should be tested rather
than answered by an impressive-looking equation.

Repetition adds a different danger. Unexamined departures can become
familiar, spread through a team, enter local routine, hide defects, or
be treated as evidence that the formal requirement was unnecessary.
Safety literature recognizes this risk of normalized
deviation.\footnote{National Aeronautics and Space Administration,
  \emph{NASA Human Factors Handbook}, NASA-HDBK-8709.25, Version 1.4
  (2023),
  https://standards.nasa.gov/sites/default/files/standards/NASA/Baseline/4/NASA-HDBK-870925-Baseline-14.pdf;
  Columbia Accident Investigation Board, \emph{Report, Volume I} (2003),
  https://ntrs.nasa.gov/citations/20030093634.} But repetition does not
mechanically lower a single internal threshold. Debriefing,
consequences, rest, training, repair, and redesigned work can interrupt
the trajectory or improve later triage.

Managers should not diagnose the phrase in workers from outside, and
workers should not diagnose it in managers merely because a decision is
harsh. What is observable is the decision and its reach. A worker's
local departure may affect one task. A decision by someone who controls
staffing, schedules, priorities, equipment, or reporting can alter the
action space of a whole unit. Power changes propagation.

A system that keeps every standard nominally mandatory after it has
removed the means to satisfy them is not preserving every standard. It
is outsourcing the choice of which one will break to the person with the
least time and often the least authority.

\hypertarget{hard-work-and-an-owned-future}{%
\subsection{Hard Work and an Owned
Future}\label{hard-work-and-an-owned-future}}

Quiet quitting proper is not anti-work. Hard effort can produce mastery,
security, ownership, portable skill, belonging, influence, or a mission
a person shares. The same person may grind for a craft, family, union
campaign, future business, or reciprocal employer and withdraw excess
from an institution where every rescue becomes the new minimum. The
objects are different.

The trouble begins when one person's fit becomes a universal character
test: \emph{I can grind here, so anyone who will not is lazy or weak.}
Sociologist Alison R. Buck uses \textbf{identity wage} to describe
identity rewards that can attract workers and secure cooperation even
alongside lower pay or stability in the games industry.\footnote{Alison
  R. Buck, ``The Price of Consent: Identity Wages in the Games
  Industry,'' \emph{Journal of Contemporary Ethnography} 51, no. 6
  (2022): 868-894, https://doi.org/10.1177/08912416221085558.} I use the
term with a related warning. Pride and toughness can be real goods. They
become suspect compensation when an institution offers the identity of
endurance in place of material or otherwise agreed reciprocity.

One influential strand of American work mythology connects the
``hard-working American'' to a secure and increasingly self-directed
future. The history is plural and its access has never been equal.
Still, the input-output story has force: disciplined effort is supposed
to improve the chance of an owned life.

Research on absolute income mobility estimates that about 90 percent of
children born in 1940 had higher inflation-adjusted household incomes
than their parents at roughly age 30, compared with about half of
children born in the 1980s.\footnote{Raj Chetty et al., ``The Fading
  American Dream: Trends in Absolute Income Mobility Since 1940,''
  \emph{Science} 356, no. 6336 (2017): 398-406,
  https://doi.org/10.1126/science.aal4617.} In a 2024 Pew survey, 53
percent of U.S. adults said the American Dream remained possible, while
41 percent said it was once possible but no longer was.\footnote{Gabriel
  Borelli, ``Americans Are Split over the State of the American Dream,''
  Pew Research Center, July 2, 2024,
  https://www.pewresearch.org/short-reads/2024/07/02/americans-are-split-over-the-state-of-the-american-dream/.}
These are population findings, not a judgment about one employer.

The decisive word is \emph{here}. A worker may still believe in effort,
family, ownership, independence, and building. What can collapse is the
belief that additional sacrifice \textbf{here} is a credible route to
those goods. No ordinary job owes limitless advancement or fulfillment.
But when an employer invokes toughness or national virtue to demand
extraordinary loyalty after the reciprocal path is gone, the bargain
becomes all additional input and no believable future output. That is
extraction wearing the language of character.

Quiet quitting proper can preserve ambition by releasing scarce capacity
for a future elsewhere. The worker is not necessarily refusing to build.
They are refusing to let one institution claim every tool with which
they might build.

\hypertarget{a-last-boundary}{%
\subsection{A Last Boundary}\label{a-last-boundary}}

The ethical test is strict because the claim is serious. The mismatch
must be persistent and material. The withdrawn effort must be
established excess rather than owed core work. Repair must be reasonably
attempted or excused by urgency, retaliation risk, incapacity, or
obvious futility. Closure must be evidenced, not merely declared. The
boundary must be proportionate. Competent performance, honesty, lawful
conduct, and professional duty must remain. Foreseeable collateral harm
must be minimized. The purpose must be truth and sustainability, not
punishment. Acute failure modes must be considered. A material repair
must reopen assessment.

The test is not so strict that the worker must consume the remainder of
their capacity proving they deserve to preserve it.

For workers, the first act can be small: name the contradiction, request
a safe priority, keep an appropriate record, and seek collective help
early where it is available. For organizers, the individual boundary is
not a substitute for building power; shared constraints are often better
met by a shared account and coordinated action. For front-line
supervisors, the task may be to validate and escalate a conflict they
cannot personally resource. For senior decision-makers, the question
after a worker's output changes should be: \emph{What unmodeled input
disappeared?}

Quiet quitting should be a last boundary, not a first weapon.

We do not manufacture failure. We stop concealing it.

We do not abandon duty. We define it.

We do not withdraw care from human beings. We withdraw unlimited
institutional access to ourselves.

\textbf{We withdraw the subsidy, not the integrity.}

\clearpage
\begingroup
\sloppy
\theendnotes
\endgroup

\end{document}
